% In this file you should put the actual content of the blueprint.
% It will be used both by the web and the print version.
% It should *not* include the \begin{document}

\part{Best-arm identification in non-stationary linear bandits}
\label{part:main}

\textbf{Overview}

This blueprint guides the Lean formalization of the paper \emph{On the Complexity of Best-Arm
Identification in Non-Stationary Linear Bandits} by Leo Maynard-Zhang, Zhihan Xiong, Kevin Jamieson
and Maryam Fazel (COLT 2026, \cite{maynardzhang2026complexity}). A learner plays arms $x_t$ from a
finite set $\X \subset \R^d$ against parameters $\theta_1, \dots, \theta_T$ fixed in advance by an
adversary, observes the rewards $x_t^\top \theta_t + \eps_t$, and must identify after $T$ rounds the
hindsight best arm $x_\star = \argmax_{x \in \X} x^\top \thetabar_T$, $\thetabar_T = \frac1T
\sum_t \theta_t$. The paper shows that the error probability of the best algorithm is
$\exp(-\Theta(T / \Hadj(\X, \Delta)))$, where
$\Hadj(\X, \Delta) = \min_\lambda \max_{(x, x') \in \I} \norm{x - x'}^2_{A(\lambda)^{-1}} / \Delta^2$
only involves the \emph{adjacent} pairs of vertices of $\conv \X$: a lower bound for every algorithm
(Theorem~\ref{thm:lower_bound}) and a matching upper bound for the algorithm Adjacent-BAI
(Theorem~\ref{thm:upper_bound}). The key geometric fact is the adjacency lemma
(Lemma~\ref{lem:adjacency}): a vertex that is not beaten by its adjacent vertices is optimal.

\textbf{Formalization.} The protocol is that of the \emph{Lean Machine Learning} library (LML):
algorithms are LML identification algorithms with a fixed budget, the environment is the
non-stationary linear bandit, and the statements quantify over all runs of the algorithms. The
statements were checked against the paper; four of them are corrected
(Remarks~\ref{rem:lower_bound_degenerate} and~\ref{rem:ls_bounds}), and several hypotheses that the
proofs do not use are dropped. The proof of the lower bound uses a construction with a factor $2$
of slack in place of the paper's half-horizon split (Remark~\ref{rem:opt_lower_route}), and the
adjacency lemma is proved from scratch by induction on the number of points, in place of the
appeal to the theory of polytopes (Chapter~\ref{chap:prereq_polytope}).

\chapter{Setting (Sections 2 and 3)}
\label{chap:setting}

Throughout, $\X \subset \R^d$ is a finite set of arms ($\R^d$ is \texttt{EuclideanSpace} over a finite type $\iota$, with
$d = \abs{\iota}$) and $T \in \N$ is the budget. Rounds are numbered $0, \dots, T - 1$.

\begin{definition}[Non-stationary linear bandit]
  \label{def:setting}
  \lean{MaynardZhang2026Complexity.avgParam, Bandits.Linear.nonstationaryLinearEnv}
  \leanok
  An instance is a parameter sequence $\theta = (\theta_t)_{t \in \N}$ in $\R^d$ and noise laws
  $(\xi_t)_t$ on $\R$: at round $t$ the learner plays $x_t \in \X$ and observes
  $r_t = x_t^\top \theta_t + \eps_t$ with $\eps_t \sim \xi_t$ independent of the past
  (\texttt{nonstationaryLinearEnv}, an LML environment whose reward kernel at round $t$ is
  the linear kernel of $\theta_t$). The noise is centered $1$-sub-Gaussian in the upper bound and
  standard Gaussian in the lower bound. The average parameter is
  $\thetabar_T = \frac1T \sum_{t < T} \theta_t$.
\end{definition}

\begin{definition}[Extreme points and adjacency, Definition 1]
  \label{def:adjacency}
  \lean{Set.vertices, Set.IsAdjacent, Set.adjacentPairs, Set.adjacentTo, Set.vertexPairs}
  \leanok
  The \emph{vertices} $\V$ of $\X$ are the exposed points of $\conv \X$. Two distinct vertices
  $x, x'$ are \emph{adjacent} if the segment $[x, x']$ is an exposed subset of $\conv \X$ (an edge).
  $\Y = \{(x, x') \in \V^2 : x \ne x'\}$ is the set of pairs of distinct vertices,
  $\I \subseteq \Y$ the set of adjacent pairs, and $\I^x$ the set of vertices adjacent to $x$.
\end{definition}

\begin{definition}[Best arm and min-gap]
  \label{def:min_gap}
  \lean{MaynardZhang2026Complexity.IsBestArm, MaynardZhang2026Complexity.minGap,
    MaynardZhang2026Complexity.MinGapGE}
  \leanok
  \uses{def:setting, def:adjacency}
  An arm $x \in \X$ is a \emph{best arm} of $\theta$ over $T$ rounds if it maximizes
  $y \mapsto y^\top \thetabar_T$ over $\X$. The \emph{min-gap} of the best arm $x_\star$ is
  $\Delta_{(1)} = \min_{x \in \V \setminus \{x_\star\}} (x_\star - x)^\top \thetabar_T$. The sequence
  $\theta$ has min-gap at least $\Delta$ if some $x \in \X$ satisfies $(x - y)^\top \thetabar_T \ge
  \Delta$ for every vertex $y \ne x$; for $\Delta > 0$ this $x$ is then the unique best arm
  (Lemma~\ref{lem:min_gap_best_arm}).
\end{definition}

\begin{definition}[Feasible sequences, Definition 2]
  \label{def:feasible}
  \lean{MaynardZhang2026Complexity.FeasibleSeqs}
  \leanok
  \uses{def:min_gap}
  Two parameter sequences are \emph{feasible} for $(\X, \Delta)$ if they have different best arms
  and each has min-gap at least $\Delta$.
\end{definition}

\begin{definition}[Designs and Mahalanobis norms]
  \label{def:design}
  \lean{Learning.IsDesign, Learning.designMatrix, Matrix.mahalanobisSq,
    MaynardZhang2026Complexity.PosDefDesign}
  \leanok
  A \emph{design} on $\X$ is a probability vector $\lambda \in \triangle_{\X}$ (a finitely supported
  function with positive weights on $\X$ summing to $1$); its design matrix is
  $A(\lambda) = \sum_{x} \lambda_x x x^\top$. For a matrix $A$, $\norm{v}_A^2 = v^\top A v$.
  \texttt{PosDefDesign} is the set of designs with positive definite design matrix.
\end{definition}

\begin{definition}[Adjacency complexity]
  \label{def:h_adjacent}
  \lean{MaynardZhang2026Complexity.adjValue, MaynardZhang2026Complexity.adjDesignValue,
    MaynardZhang2026Complexity.hAdjacent}
  \leanok
  \uses{def:adjacency, def:design}
  \[
    \Hadj(\X, \Delta) = \frac{1}{\Delta^2} \min_{\lambda} \max_{(x, x') \in \I}
      \norm{x - x'}^2_{A(\lambda)^{-1}},
  \]
  the minimum being over the designs with positive definite design matrix (an infimum in Lean,
  \texttt{adjDesignValue}; the inner maximum is \texttt{adjValue}).
\end{definition}

\begin{lemma}[Adjacency lemma, Lemma 1]
  \label{lem:adjacency}
  \lean{MaynardZhang2026Complexity.exists_pos_inner_iff}
  \leanok
  \uses{def:adjacency, lem:local_global}
  Let $x \in \V$ and $\theta \in \R^d$. There is $y \in \X$ with $(y - x)^\top \theta > 0$ if and
  only if there is $z \in \I^x$ with $(z - x)^\top \theta > 0$.
\end{lemma}

\begin{proof}
  \leanok
  \uses{lem:vertices_subset, lem:local_global}
  Backward: adjacent vertices are points of $\X$ (Lemma~\ref{lem:vertices_subset}). Forward:
  Lemma~\ref{lem:local_global}.
\end{proof}

\begin{remark}
  The paper deduces Lemma~\ref{lem:adjacency} from the cone property of polytopes,
  $\conv \X \subseteq x + \cone\{z - x : z \in \I^x\}$ (\cite[Lemma 3.6]{ziegler1995lectures}),
  which is not in Mathlib. The blueprint proves the lemma directly
  (Lemma~\ref{lem:local_global}) and deduces the cone property in the form needed for
  Lemma~\ref{lem:ratio_adjacent} (Lemma~\ref{lem:cone_form}).
\end{remark}

\begin{lemma}[The arm of the min-gap is the best arm]
  \label{lem:min_gap_best_arm}
  \lean{Set.Finite.mem_vertices_of_forall_vertices}
  \leanok
  \uses{def:min_gap, lem:strict_expose, lem:max_face}
  Let $\Delta > 0$ and let $x \in \X$ beat every other vertex by $\Delta$ for $\thetabar_T$. Then $x$
  is a vertex, $x^\top \thetabar_T > y^\top \thetabar_T$ for all $y \in \X \setminus \{x\}$, and $x$
  is the unique best arm.
\end{lemma}

\begin{proof}
  \leanok
  \uses{lem:strict_expose, lem:max_face}
  Every vertex $v$ has $v^\top \thetabar_T \le x^\top \thetabar_T$, with equality only for $v = x$.
  Let $y \in \X$ with $y^\top \thetabar_T \ge x^\top \thetabar_T$. As $y \in \conv \X = \conv \V$
  (Lemma~\ref{lem:strict_expose}), $y$ is a convex combination of the vertices $v$ with
  $v^\top \thetabar_T = x^\top \thetabar_T$ (Lemma~\ref{lem:max_face}), of which there is at most
  one, $x$: so this set is $\{x\}$, $x$ is a vertex and $y = x$. Applied to $y = x$ this shows
  that $x$ is a vertex, and applied to $y \ne x$ that $y^\top \thetabar_T < x^\top \thetabar_T$.
\end{proof}

\chapter{Lower bound (Section 4 and Appendix A)}
\label{chap:lower_bound}

\section{A lower bound for non-stationary bandits}

A non-stationary bandit with finitely many arms $\mathcal{A}$ is a sequence
$\nu = (\nu_{t, a})_{t \in \N, a \in \mathcal{A}}$ of probability measures on $\R$, played through
LML's environment \texttt{Environment.banditSeq}. For an identification algorithm with budget
$T$, $X_t$ is the arm of round $t$ and $\hat k$ the output.

\begin{lemma}[Non-stationary Kaufmann--Capp\'e--Garivier bound, Lemma 2]
  \label{lem:gen_lower}
  \lean{MaynardZhang2026Complexity.le_max_prob_ne}
  \leanok
  \uses{lem:bretagnolle_huber, lem:data_processing_output, lem:chain_rule_nonstationary}
  Let $\nu, \nu'$ be non-stationary bandits with $\KL(\nu_{t, a} \,\|\, \nu'_{t, a}) < \infty$ for
  all $t, a$, let $k \ne k'$ be arms and let $A$ be an identification algorithm with budget $T$.
  For every run of $A$ against $\nu$ (law $\Prob_\nu$) and every run against $\nu'$,
  \[
    \max\left(\Prob_\nu(\hat k \ne k), \Prob_{\nu'}(\hat k \ne k')\right)
    \ge \frac14 \exp\left(-\sum_{a} \sum_{t < T} \Prob_\nu(X_t = a)\,
      \KL(\nu_{t, a} \,\|\, \nu'_{t, a})\right).
  \]
\end{lemma}

\begin{proof}
  \leanok
  \uses{lem:bretagnolle_huber, lem:data_processing_output, lem:chain_rule_nonstationary}
  Let $E = \{\hat k \ne k\}$; since $k \ne k'$, $E^c \subseteq \{\hat k \ne k'\}$. By the
  Bretagnolle--Huber inequality (Lemma~\ref{lem:bretagnolle_huber}) for the laws of the outputs,
  $\Prob_\nu(E) + \Prob_{\nu'}(E^c) \ge \frac12 \exp(-\KL)$, so the maximum is at least
  $\frac14 \exp(-\KL)$, where $\KL$ is the divergence between the laws of the outputs. It is at
  most the divergence between the laws of the histories of the $T$ rounds
  (Lemma~\ref{lem:data_processing_output}), which is the sum in the statement
  (Lemma~\ref{lem:chain_rule_nonstationary}).
\end{proof}

\begin{remark}
  The paper states Lemma~\ref{lem:gen_lower} for the best arms $k$ and $k'$ of $\nu$ and $\nu'$,
  assumed different; only $k \ne k'$ is used. Its proof goes through the change of measure at a
  stopping time (Lemma~\ref{lem:change_of_measure}) and Lemma 19 of
  \cite{kaufmann2016complexity}; the route above, through the chain rule for the divergence of the
  histories, avoids both.
\end{remark}

\begin{lemma}[Change of measure at a stopping time, Lemma 6]
  \label{lem:change_of_measure}
  \lean{MaynardZhang2026Complexity.trajMeasure_eq_lintegral_exp_neg_llrProcess,
    Bandits.llrProcess,
    Bandits.trajMeasure_banditSeq_apply_eq_setLIntegral_of_isStoppingTime,
    Bandits.trajMeasure_banditSeq_apply_eq_setLIntegral,
    Bandits.feedback_banditSeq_eq_withDensity,
    Learning.IsAlgEnvSeq.hasLaw_history_withDensity_of_feedback_eq,
    Learning.stepKernel_eq_withDensity_of_feedback_eq}
  \leanok
  Let $\nu, \nu'$ be non-stationary bandits with countably many arms and $\nu_{t,a}$, $\nu'_{t,a}$
  mutually absolutely continuous, and let $\Prob_\nu$, $\Prob_{\nu'}$ be the laws of the
  trajectory of an algorithm against them. Let
  $L_n = \sum_{s < n} \log \frac{d\nu_{s, X_s}}{d\nu'_{s, X_s}}(Y_s)$ be the log-likelihood ratio
  of the first $n$ rounds, and $\mathcal{F}_n$ the $\sigma$-algebra generated by the rounds
  $0, \dots, n$. For every stopping time $\sigma$ of this filtration, finite everywhere, and every
  $E \in \mathcal{F}_\sigma$, $\Prob_{\nu'}(E) = \E_\nu[\indic_E \exp(-L_{\sigma + 1})]$.
\end{lemma}

\begin{proof}
  \leanok
  The step kernel of round $n$ against $\nu'$ has density
  $(h, (a, y)) \mapsto \frac{d\nu'_{n, a}}{d\nu_{n, a}}(y)
  = \exp(-\log \frac{d\nu_{n, a}}{d\nu'_{n, a}}(y))$ ($\nu_{n, a}$-almost everywhere, by mutual
  absolute continuity) with respect to that against $\nu$, so by induction on $n$ the law of the
  history of the first $n$ rounds against $\nu'$ has density $\exp(-L_n)$ with respect to that
  against $\nu$. This is the identity for $E \in \mathcal{F}_n$ (determined by the first $n + 1$
  rounds) with $L_{n + 1}$. For $E \in \mathcal{F}_\sigma$, $E$ is the disjoint union of the sets
  $E \cap \{\sigma = n\} \in \mathcal{F}_n$, on which $L_{\sigma + 1} = L_{n + 1}$.
\end{proof}

\begin{remark}
  \label{rem:change_of_measure_index}
  In the paper the rounds are numbered from $1$, and $\mathcal{F}_t$ and $L_t$ cover the same
  rounds $1, \dots, t$: its $L_\sigma$ is $L_{\sigma + 1}$ here. The statement written in phase 1
  of this formalization, with $L_\sigma$, was false (an indexing error of the translation, not of
  the paper): for the constant stopping time $\sigma = 0$ and the event
  $E = \{Y_0 \in B\} \in \mathcal{F}_0$, its left-hand side is $\E[\nu'_{0, X_0}(B)]$ and its
  right-hand side $\Prob_\nu(E) = \E[\nu_{0, X_0}(B)]$.
\end{remark}

\section{The optimization problem}

\begin{definition}[Optimization problem of Lemma 3]
  \label{def:opt_value}
  \lean{MaynardZhang2026Complexity.PairFeasible, MaynardZhang2026Complexity.pairValue,
    MaynardZhang2026Complexity.optValue}
  \leanok
  \uses{def:adjacency, def:design}
  For a pair $(x, x')$ of vertices, $(\theta, v)$ is \emph{feasible} if
  $(x - y)^\top \theta \ge \Delta$ for every vertex $y \ne x$ and $(x' - y)^\top (\theta + v) \ge
  \Delta$ for every vertex $y \ne x'$. For a matrix $A$, $p_A(x, x') = \inf\{v^\top A v : (\theta,
  v) \text{ feasible}\}$, and
  \[
    f(\X, \Delta) = \max_{\lambda \in \triangle_{\X}} \min_{(x, x') \in \Y} p_{A(\lambda)}(x, x'),
  \]
  the maximum being over all designs (a supremum in Lean).
\end{definition}

\begin{lemma}[Feasible points exist]
  \label{lem:pair_feasible_nonempty}
  \lean{MaynardZhang2026Complexity.exists_pairFeasible,
    MaynardZhang2026Complexity.exists_pairFeasible_of_finite,
    MaynardZhang2026Complexity.exists_forall_le_inner_sub_of_mem_vertices,
    MaynardZhang2026Complexity.pairValue_nonneg}
  \leanok
  \uses{def:opt_value, lem:strict_expose}
  For distinct vertices $x, x'$ there is a feasible $(\theta, v)$. For $A$ positive semidefinite,
  $p_A(x, x') \ge 0$.
\end{lemma}

\begin{proof}
  \leanok
  \uses{lem:strict_expose}
  Let $w$, $w'$ strictly expose $x$ and $x'$ (Lemma~\ref{lem:strict_expose}); for $c$ large,
  $\theta = c w$ and $\theta + v = c w'$ satisfy the finitely many constraints. The values
  $v^\top A v$ are nonnegative.
\end{proof}

\begin{lemma}[Cauchy--Schwarz inequality for Mahalanobis norms]
  \label{lem:mahalanobis_cauchy_schwarz}
  \lean{Matrix.PosDef.inner_sq_le_mahalanobisSq_inv_mul}
  \leanok
  \uses{def:design}
  For $A$ positive definite and $z, v \in \R^d$, $(z^\top v)^2 \le \norm{z}^2_{A^{-1}}
  \norm{v}^2_A$.
\end{lemma}

\begin{proof}
  \leanok
  $z^\top v = (A^{-1/2} z)^\top (A^{1/2} v)$, and $\norm{A^{-1/2} z}^2 = \norm{z}^2_{A^{-1}}$,
  $\norm{A^{1/2} v}^2 = \norm{v}^2_A$; apply the Cauchy--Schwarz inequality of $\R^d$.
\end{proof}

\begin{lemma}[Lower bound on the inner problem, Lemma 7]
  \label{lem:pair_value_lower}
  \lean{MaynardZhang2026Complexity.le_pairValue,
    MaynardZhang2026Complexity.PairFeasible.two_mul_le_inner,
    MaynardZhang2026Complexity.PairFeasible.div_mahalanobisSq_inv_le}
  \leanok
  \uses{def:opt_value, lem:pair_feasible_nonempty, lem:mahalanobis_cauchy_schwarz}
  Let $\Delta > 0$, $A$ positive definite and $(x, x') \in \Y$. Then
  $p_A(x, x') \ge 4\Delta^2 / \norm{x - x'}^2_{A^{-1}}$.
\end{lemma}

\begin{proof}
  \leanok
  \uses{lem:pair_feasible_nonempty, lem:mahalanobis_cauchy_schwarz}
  The constraints at $y = x'$ and at $y = x$ add up to $(x' - x)^\top v \ge 2\Delta$ for every
  feasible $(\theta, v)$. By the Cauchy--Schwarz inequality for the inner products of $A^{-1}$ and
  $A$ (Lemma~\ref{lem:mahalanobis_cauchy_schwarz}),
  $4 \Delta^2 \le ((x' - x)^\top v)^2 \le \norm{x - x'}^2_{A^{-1}} \norm{v}^2_A$. The infimum is
  over a nonempty set (Lemma~\ref{lem:pair_feasible_nonempty}).
\end{proof}

\begin{lemma}[Upper bound on the inner problem, Lemma 8]
  \label{lem:pair_value_upper}
  \lean{MaynardZhang2026Complexity.pairValue_le,
    MaynardZhang2026Complexity.exists_pairFeasible_mahalanobisSq_eq}
  \leanok
  \uses{def:opt_value, lem:is_adjacent_iff, lem:pair_feasible_nonempty}
  Let $A$ be positive definite and $(x, x') \in \I$. Then
  $p_A(x, x') \le 4\Delta^2 / \norm{x - x'}^2_{A^{-1}}$.
\end{lemma}

\begin{proof}
  \leanok
  \uses{lem:is_adjacent_iff, lem:pair_feasible_nonempty}
  Let $w$ and $\eps$ be given by Lemma~\ref{lem:is_adjacent_iff}, $c = \norm{x - x'}^2_{A^{-1}} > 0$,
  $\theta = \frac{\Delta}{c} A^{-1}(x - x') + \alpha w$ and $v = -\frac{2\Delta}{c} A^{-1}(x - x')$.
  Then $v^\top A v = 4\Delta^2 / c$, the constraints at $y = x'$ and $y = x$ hold with equality
  ($(x - x')^\top w = 0$), and the other constraints hold for $\alpha$ large since
  $(x - y)^\top w = (x' - y)^\top w \ge \eps$. The values are nonnegative, so the infimum is at most
  $4\Delta^2 / c$.
\end{proof}

\begin{remark}
  The paper assumes $\Delta > 0$ in Lemma~\ref{lem:pair_value_upper}; the construction works for
  every $\Delta$. The same construction with $\theta = -v/2 + \alpha w$ shows that, for an adjacent
  pair, the feasible $v$ are exactly the half-space $(x' - x)^\top v \ge 2\Delta$.
\end{remark}

\begin{lemma}[Closed form for adjacent pairs, Lemma 4]
  \label{lem:equality}
  \lean{MaynardZhang2026Complexity.pairValue_eq}
  \leanok
  \uses{lem:pair_value_lower, lem:pair_value_upper}
  Let $\Delta > 0$, $A$ positive definite and $(x, x') \in \I$. Then
  $p_A(x, x') = 4\Delta^2 / \norm{x - x'}^2_{A^{-1}}$.
\end{lemma}

\begin{proof}
  \leanok
  \uses{lem:pair_value_lower, lem:pair_value_upper}
  Lemmas~\ref{lem:pair_value_lower} and~\ref{lem:pair_value_upper}. (The paper states it for the
  design matrices of full rank.)
\end{proof}

\begin{lemma}[Monotonicity of the inner problem]
  \label{lem:pair_value_mono}
  \lean{MaynardZhang2026Complexity.pairValue_mono, MaynardZhang2026Complexity.pairValue_smul}
  \leanok
  \uses{def:opt_value, lem:pair_feasible_nonempty}
  For every pair $(x, x')$, $A \mapsto p_A(x, x')$ is monotone for the Loewner order on positive
  semidefinite matrices, and $p_{cA}(x, x') = c\, p_A(x, x')$ for $c \ge 0$.
\end{lemma}

\begin{proof}
  \leanok
  \uses{lem:pair_feasible_nonempty}
  Pointwise in $v$ ($A \le B$ gives $v^\top A v \le v^\top B v$); the feasible set does not depend
  on $A$, and the values are nonnegative (Lemma~\ref{lem:pair_feasible_nonempty}). (In Lean, the
  infimum over an empty feasible set is $0$ for every $A$.)
\end{proof}

\begin{lemma}[Adjacent pairs exist]
  \label{lem:exists_adjacent_pair}
  \lean{Set.mem_vertices_of_forall_norm_le, Set.Finite.exists_isAdjacent,
    MaynardZhang2026Complexity.nonempty_vertexPairs}
  \leanok
  \uses{lem:strict_expose, lem:local_global}
  A point of $\X$ of maximal norm is a vertex. If $\X$ has at least two points, it has two
  adjacent vertices; in particular $\Y \ne \emptyset$.
\end{lemma}

\begin{proof}
  \leanok
  \uses{lem:strict_expose, lem:local_global}
  If $\norm{x}$ is maximal and $y \in \X \setminus \{x\}$, then
  $0 < \norm{y - x}^2 \le 2 \norm{x}^2 - 2 y^\top x$, so $y^\top x < x^\top x$: $x$ is strictly
  exposed by itself (Lemma~\ref{lem:strict_expose}). Lemma~\ref{lem:local_global} gives a vertex
  adjacent to it.
\end{proof}

\begin{lemma}[Positivity of $f$]
  \label{lem:opt_value_pos}
  \lean{MaynardZhang2026Complexity.optValue_pos,
    MaynardZhang2026Complexity.exists_pairFeasible_lt_two_mul_optValue,
    MaynardZhang2026Complexity.bddAbove_range_iInf_pairValue,
    MaynardZhang2026Complexity.iInf_pairValue_le_optValue,
    MaynardZhang2026Complexity.le_pairValue_inv_card_smul_sum_outerSelf}
  \leanok
  \uses{def:opt_value, lem:pair_feasible_nonempty}
  If $\Delta > 0$ and $\X$ has at least two points, then $0 < f(\X, \Delta) < \infty$ (the supremum
  over designs is bounded). Consequently, for every design $\lambda$ some $(x, x') \in \Y$ has a
  feasible $(\theta, v)$ with $v^\top A(\lambda) v < 2 f(\X, \Delta)$.
\end{lemma}

\begin{proof}
  \leanok
  \uses{lem:pair_feasible_nonempty, lem:exists_adjacent_pair, lem:pair_value_lower}
  $\Y$ is nonempty (Lemma~\ref{lem:exists_adjacent_pair}) and finite. Bounded: fix a pair of $\Y$
  and a feasible $(\theta, v)$ for it; for every design,
  $p_{A(\lambda)} \le v^\top A(\lambda) v = \sum_x \lambda_x (x^\top v)^2 \le \sum_{y \in \X}
  (y^\top v)^2$. Positive: let $A_u = \frac1n \sum_{y \in \X} y y^\top$ be the design matrix of the
  uniform design, $n = \abs{\X}$. For $(x, x') \in \Y$ every feasible $v$ has
  $2\Delta \le x'^\top v - x^\top v$ (first step of Lemma~\ref{lem:pair_value_lower}), so
  $4 \Delta^2 \le 2 ((x^\top v)^2 + (x'^\top v)^2) \le 2 n\, v^\top A_u v$ (as $x, x' \in \X$), and
  $p_{A_u}(x, x') \ge 2\Delta^2 / n$; take the minimum over the finitely many pairs. Finally,
  $\min_{\Y} p_{A(\lambda)} \le f(\X, \Delta) < 2 f(\X, \Delta)$: some pair has
  $p_{A(\lambda)}(x, x') < 2 f(\X, \Delta)$, and some feasible point of that pair is below
  $2 f(\X, \Delta)$.
\end{proof}

\section{Lower bounds on the error}

\begin{lemma}[Runs agreeing before a round]
  \label{lem:actions_law_agree}
  \lean{Learning.IsAlgEnvSeq.identDistrib_action_of_env_eq,
    Learning.IsAlgEnvSeq.identDistrib_history_of_env_eq, Learning.IsAlgEnvSeqUntil.of_env_eq}
  \leanok
  Let two environments have the same feedback kernels at rounds $0, \dots, t - 1$. For a given
  algorithm, the law of $X_t$ (and of the history of the first $t$ rounds) is the same in every run
  against either environment.
\end{lemma}

\begin{proof}
  \leanok
  The law of the history of the first $t$ rounds followed by the action of round $t$ is the
  composition of the policy kernels of rounds $0, \dots, t$ and of the feedback kernels of rounds
  $0, \dots, t - 1$ (uniqueness of the law of a run).
\end{proof}

\begin{lemma}[Optimization-based lower bound, Lemma 3]
  \label{lem:opt_lower}
  \lean{MaynardZhang2026Complexity.exists_minGapGE_le_max_error_optValue,
    MaynardZhang2026Complexity.isBestArm_iff_eq_of_forall_vertices,
    MaynardZhang2026Complexity.avgParam_ite_eq, MaynardZhang2026Complexity.avgParam_add_const,
    Learning.designOfWeights, Learning.isDesign_designOfWeights}
  \leanok
  \uses{lem:gen_lower, def:opt_value, def:min_gap, lem:opt_value_pos, lem:actions_law_agree,
    lem:kl_gaussian, lem:min_gap_best_arm}
  Let $\X$ have at least two points, $\Delta > 0$, $T \ge 1$ and $A$ an identification algorithm
  with budget $T$. There are parameter sequences $\theta$, $\theta'$, both with min-gap at least
  $\Delta$, such that for every run of $A$ against $\theta$ and against $\theta'$ with standard
  Gaussian noise,
  \[
    \max\left(\Prob_\theta(\hat x \text{ is not a best arm}),
      \Prob_{\theta'}(\hat x \text{ is not a best arm})\right)
    \ge \frac14 \exp\left(-T f(\X, \Delta)\right).
  \]
\end{lemma}

\begin{proof}
  \leanok
  \uses{lem:gen_lower, lem:opt_value_pos, lem:actions_law_agree, lem:kl_gaussian,
    lem:min_gap_best_arm}
  Let $P_0$ be the law of a run of $A$ against $\theta \equiv 0$ and
  $\lambda_x = \frac1T \sum_{t < T} P_0(X_t = x)$, a design. Let $(x, x') \in \Y$ minimize
  $p_{A(\lambda)}$; since $p_{A(\lambda)}(x, x') \le f(\X, \Delta) < 2 f(\X, \Delta)$
  (Lemma~\ref{lem:opt_value_pos}), there is a feasible $(\tilde\theta, \tilde v)$ with
  $\tilde v^\top A(\lambda) \tilde v \le 2 f(\X, \Delta)$. Let
  \[
    \theta_t = 0,\ \theta'_t = \tilde v \ (t < T - 1), \qquad
    \theta_{T - 1} = T \tilde\theta,\ \theta'_{T - 1} = T \tilde\theta + \tilde v.
  \]
  Then $\thetabar_T = \tilde\theta$ and $\thetabar'_T = \tilde\theta + \tilde v$, so by feasibility
  both sequences have min-gap at least $\Delta$, with best arms $x$ and $x'$
  (Lemma~\ref{lem:min_gap_best_arm}). The law of $X_t$, $t \le T - 1$, under $\theta$ depends only
  on $\theta_0, \dots, \theta_{T - 2} = 0$, so it is its law under $P_0$
  (Lemma~\ref{lem:actions_law_agree}). The divergence between the reward laws of arm $a$ at round
  $t$ is $(a^\top \tilde v)^2 / 2$ (Lemma~\ref{lem:kl_gaussian}), so Lemma~\ref{lem:gen_lower} with
  $k = x$, $k' = x'$ gives the bound with exponent
  $\sum_{t < T} \sum_a P_0(X_t = a) (a^\top \tilde v)^2 / 2 = \frac{T}{2} \tilde v^\top A(\lambda)
  \tilde v \le T f(\X, \Delta)$. The error events $\{\hat x \ne x\}$ and $\{\hat x \ne x'\}$ are the
  events that $\hat x$ is not a best arm, the best arms being unique.
\end{proof}

\begin{remark}
  \label{rem:opt_lower_route}
  The paper takes $\theta_t = 0$, $\theta'_t = v$ in the first half of the horizon and
  $\theta_t = \theta'_t = \theta$ in the second, ``assuming without loss of generality that $T$ is
  even''. This yields the exponent $T \tilde v^\top A(\lambda) \tilde v$, so the bound
  $\exp(-T f)$ needs a feasible point attaining the inner infimum $p_{A(\lambda)}$ (true, by a
  Frank--Wolfe argument for a convex quadratic on a polyhedron, but not available in Mathlib), and
  for $T = 1$ the first half is the whole horizon, where $\thetabar = 0$ has no min-gap. The
  construction above spreads the difference between the instances over all $T$ rounds, while the
  parameters of $\theta$ before the last round vanish, so that the law of the actions under
  $\theta$ is fixed before $(\tilde\theta, \tilde v)$ is chosen. This halves the exponent, so no
  attainment is needed, and works for every $T \ge 1$.
\end{remark}

\begin{lemma}[Bound on $f$]
  \label{lem:opt_value_le}
  \lean{MaynardZhang2026Complexity.optValue_le, MaynardZhang2026Complexity.hAdjacent_pos,
    MaynardZhang2026Complexity.adjDesignValue_pos,
    MaynardZhang2026Complexity.iInf_pairValue_le_one_add_mul}
  \leanok
  \uses{lem:equality, lem:pair_value_mono, def:h_adjacent, lem:exists_adjacent_pair,
    lem:mahalanobis_cauchy_schwarz}
  If $\X$ spans $\R^d$ and has at least two points and $\Delta > 0$, then the value of the adjacent
  design is positive and $f(\X, \Delta) \le 4 / \Hadj(\X, \Delta)$.
\end{lemma}

\begin{proof}
  \leanok
  \uses{lem:equality, lem:pair_value_mono, lem:exists_adjacent_pair,
    lem:mahalanobis_cauchy_schwarz}
  Since $\X$ spans, there is a positive definite design matrix $A_0$ (the uniform design on a basis
  contained in $\X$). Positivity: let $(x, x') \in \I$ (Lemma~\ref{lem:exists_adjacent_pair}),
  $z = x - x' \ne 0$ and $K = \sum_{y \in \X} (y^\top z)^2$. For every design $\lambda$ with
  $A(\lambda)$ positive definite, Lemma~\ref{lem:mahalanobis_cauchy_schwarz} and
  $z^\top A(\lambda) z = \sum_y \lambda_y (y^\top z)^2 \le K$ give
  $\norm{z}^4 \le \norm{z}^2_{A(\lambda)^{-1}} K$, so (applied to $A_0$) $K > 0$ and the value of
  $\lambda$ is at least $\norm{z}^2_{A(\lambda)^{-1}} \ge \norm{z}^4 / K > 0$.
  Let $a$ be the value of the adjacent design, $\lambda$ a design and $\eps > 0$. Then
  $A_\eps = (A(\lambda) + \eps A_0)/(1 + \eps)$ is the positive definite design matrix of a design,
  and $A(\lambda) \le (1 + \eps) A_\eps$. Let $(x, x') \in \I$ maximize
  $\norm{x - x'}^2_{A_\eps^{-1}}$, so that $\norm{x - x'}^2_{A_\eps^{-1}} \ge a$. By
  Lemmas~\ref{lem:pair_value_mono} and~\ref{lem:equality},
  \[
    \min_{\Y} p_{A(\lambda)} \le p_{A(\lambda)}(x, x') \le (1 + \eps)\, p_{A_\eps}(x, x')
    = (1 + \eps) \frac{4 \Delta^2}{\norm{x - x'}^2_{A_\eps^{-1}}}
    \le (1 + \eps) \frac{4\Delta^2}{a}.
  \]
  Let $\eps \to 0$ and take the supremum over $\lambda$: $f(\X, \Delta) \le 4\Delta^2 / a = 4 /
  \Hadj(\X, \Delta)$.
\end{proof}

\begin{theorem}[Lower bound, Theorem 1]
  \label{thm:lower_bound}
  \lean{MaynardZhang2026Complexity.exists_minGapGE_le_max_error}
  \leanok
  \uses{lem:opt_lower, lem:opt_value_le}
  Let $\X$ span $\R^d$ and have at least two points, $\Delta > 0$, $T \ge 1$ and $A$ an
  identification algorithm with budget $T$. There are parameter sequences $\theta$, $\theta'$, both
  with min-gap at least $\Delta$, such that for every run of $A$ against $\theta$ and against
  $\theta'$ with standard Gaussian noise,
  \[
    \max\left(\Prob_\theta(\hat x \text{ is not a best arm}),
      \Prob_{\theta'}(\hat x \text{ is not a best arm})\right)
    \ge \frac14 \exp\left(-\frac{4T}{\Hadj(\X, \Delta)}\right).
  \]
\end{theorem}

\begin{proof}
  \leanok
  \uses{lem:opt_lower, lem:opt_value_le}
  Lemma~\ref{lem:opt_lower} and $T f(\X, \Delta) \le 4T / \Hadj(\X, \Delta)$
  (Lemma~\ref{lem:opt_value_le}).
\end{proof}

\begin{remark}
  \label{rem:lower_bound_degenerate}
  The paper states Theorem~\ref{thm:lower_bound} and Lemma~\ref{lem:opt_lower} without $T \ge 1$
  and without requiring two arms. Both are needed: for $T = 0$, $\thetabar_0 = 0$ and no parameter
  sequence has min-gap $\Delta > 0$; for $\X = \{x\}$ there are no adjacent pairs, so
  $\Hadj(\X, \Delta) = 0$ and $f(\X, \Delta) = 0$ (with the conventions $a / 0 = 0$ and
  $\min \emptyset = 0$), and the bound would claim an error of at least $1/4$, while every algorithm
  outputs the only arm, which is a best arm. Lemma~\ref{lem:opt_lower} holds without the assumption
  that $\X$ spans $\R^d$.
\end{remark}

\chapter{Matching upper bound (Section 5 and Appendix B)}
\label{chap:upper_bound}

\begin{definition}[Adjacent-BAI, Algorithm 1]
  \label{def:adjacent_bai}
  \lean{MaynardZhang2026Complexity.allocMatrix, MaynardZhang2026Complexity.IsAdjacentRounding,
    MaynardZhang2026Complexity.lsEstimator, MaynardZhang2026Complexity.recommend,
    MaynardZhang2026Complexity.adjacentBAI, Learning.Algorithm.randomOrder,
    Learning.IdentAlg.fixedBudget}
  \leanok
  \uses{def:h_adjacent, def:setting}
  Let $x_0, \dots, x_{T-1} \in \X$ be an allocation, $A = \sum_{t} x_t x_t^\top$. It satisfies the
  \emph{rounding guarantee} if $\frac1T A$ is positive definite and
  $\max_{(x, x') \in \I} \norm{x - x'}^2_{(A/T)^{-1}} \le 2 \min_\lambda \max_{(x, x') \in \I}
  \norm{x - x'}^2_{A(\lambda)^{-1}}$ (Eq.~(roundingerror); the rounding procedure of
  \cite{pukelsheim2006optimal} provides such an allocation when $T \ge d^2$). Adjacent-BAI plays the
  allocation in uniformly random order (at each round, a uniform arm among those of the allocation
  not yet played), computes the least-squares estimator
  $\thetahat_T = A^{-1} \sum_{t < T} r_t X_t$ from the arms $X_t$ played and the rewards $r_t$, and
  outputs $\argmax_{x \in \X} x^\top \thetahat_T$ (with a fixed tie-breaking).
\end{definition}

\begin{lemma}[Sum over a random permutation, Lemma 9]
  \label{lem:perm_subgaussian}
  \lean{MaynardZhang2026Complexity.hasSubgaussianMGF_sum_perm}
  \leanok
  \uses{lem:sampling_without_replacement, lem:column_bessel}
  Let $x_0, \dots, x_{T-1}$ span $\R^d$ with $\norm{x_t} \le B$, $\theta_0, \dots, \theta_{T-1}$ with
  $\norm{\theta_t} \le M$, $\bar\theta = \frac1T \sum_t \theta_t$, $A = \sum_t x_t x_t^\top$ and $\pi$
  a uniform permutation of $\{0, \dots, T - 1\}$. For every $z$,
  $S = \sum_t (z^\top A^{-1} x_t) \, x_t^\top (\theta_{\pi(t)} - \bar\theta)$ is sub-Gaussian with
  variance proxy $8 M^2 B^2 \norm{z}^2_{A^{-1}}$.
\end{lemma}

\begin{proof}
  \uses{lem:sampling_without_replacement, lem:column_bessel}
  $S = \sum_t a_t^\top (\theta_{\pi(t)} - \bar\theta)$ with $a_t = x_t x_t^\top A^{-1} z$. By
  Lemma~\ref{lem:sampling_without_replacement}, $S$ is sub-Gaussian with proxy
  $4M^2 \sum_{t < T - 1} \norm{d_t}^2$, $d_t = a_t - \frac{1}{T - 1 - t} \sum_{s > t} a_s$. By
  Lemma~\ref{lem:column_bessel}, $\sum_t \norm{d_t}^2 \le 2 \sum_t \norm{a_t}^2$, and
  $\norm{a_t} = \abs{x_t^\top A^{-1} z} \norm{x_t} \le B \abs{x_t^\top A^{-1} z}$, with
  $\sum_t (x_t^\top A^{-1} z)^2 = z^\top A^{-1} A A^{-1} z = \norm{z}^2_{A^{-1}}$.
\end{proof}

\begin{remark}
  The paper proves Lemma~\ref{lem:perm_subgaussian} with the Doob martingale of $S$ for the
  filtration of $\pi(0), \dots, \pi(t)$ and an Azuma--Hoeffding argument; the induction of
  Lemma~\ref{lem:sampling_without_replacement} is the same computation without the martingale
  machinery.
\end{remark}

\begin{lemma}[Runs of Adjacent-BAI]
  \label{lem:run_adjacent_bai}
  \uses{def:adjacent_bai}
  In a run of Adjacent-BAI, there is a uniform permutation $\pi$ such that $X_t = x_{\pi(t)}$ for
  $t < T$ almost surely, and the noises $\eps_t = r_t - X_t^\top \theta_t$, $t < T$, are independent
  with laws $\xi_t$ and independent of $\pi$.
\end{lemma}

\begin{proof}
  Build the random variables on an explicit probability space ($\pi$ uniform, independent noises),
  check that they form a run (the law of the next arm given the past is uniform on the arms not yet
  played, which is the law of $x_{\pi(t)}$ given $\pi(0), \dots, \pi(t-1)$), and use the uniqueness
  of the law of a run.
\end{proof}

\begin{lemma}[Sub-Gaussian error of the least-squares estimator, Lemma 5]
  \label{lem:ls_subgaussian}
  \lean{MaynardZhang2026Complexity.hasSubgaussianMGF_lsEstimator}
  \leanok
  \uses{lem:perm_subgaussian, lem:run_adjacent_bai, lem:weighted_noise_subgaussian}
  Let the allocation have $A = \sum_t x_t x_t^\top$ positive definite and $\norm{x_t} \le 1$, let
  $\norm{\theta_t} \le 1$ for $t < T$, and let the noise be centered $1$-sub-Gaussian. In every run of
  Adjacent-BAI, for every $z$, $z^\top(\thetahat_T - \thetabar_T)$ is sub-Gaussian with variance
  proxy $9 \norm{z}^2_{A^{-1}}$.
\end{lemma}

\begin{proof}
  \uses{lem:perm_subgaussian, lem:run_adjacent_bai, lem:weighted_noise_subgaussian}
  By Lemma~\ref{lem:run_adjacent_bai}, $r_t = x_{\pi(t)}^\top \theta_t + \eps_t$, and since
  $\sum_s x_s x_s^\top \thetabar_T = A \thetabar_T$,
  \[
    z^\top(\thetahat_T - \thetabar_T)
    = \sum_s (z^\top A^{-1} x_s)\, x_s^\top (\theta_{\pi^{-1}(s)} - \thetabar_T)
      + \sum_t (z^\top A^{-1} X_t)\, \eps_t = S + \eta.
  \]
  $S$ is the statistic of Lemma~\ref{lem:perm_subgaussian} for the uniform permutation $\pi^{-1}$,
  with $B = M = 1$: sub-Gaussian with proxy $8 \norm{z}^2_{A^{-1}}$. Conditionally on $\pi$, $\eta$
  is a weighted sum of independent $1$-sub-Gaussian variables with weights
  $z^\top A^{-1} x_{\pi(t)}$, sub-Gaussian with proxy $\sum_s (z^\top A^{-1} x_s)^2 =
  \norm{z}^2_{A^{-1}}$ (Lemma~\ref{lem:weighted_noise_subgaussian}). Hence
  $\E[e^{\mu(S + \eta)}] = \E[e^{\mu S} \E[e^{\mu \eta} \mid \pi]] \le e^{\mu^2 (8 + 1)
  \norm{z}^2_{A^{-1}} / 2}$.
\end{proof}

\begin{remark}
  \label{rem:ls_bounds}
  The paper states Lemma~\ref{lem:ls_subgaussian} without the bounds $\norm{x_t} \le 1$ and
  $\norm{\theta_t} \le 1$, which are the assumptions of Theorem~\ref{thm:upper_bound} and are used
  by its proof (through Lemma~\ref{lem:perm_subgaussian} with $B = M = 1$). Without them the lemma
  is false: for $d = 1$, $T = 2$, $x = (1, 2)$, $\theta = (-10, 10)$, the term $S$ is
  $\pm 6 z$ with probability $1/2$ each, of variance $36 z^2$, while
  $9 \norm{z}^2_{A^{-1}} = 9 z^2 / 5$ (a sub-Gaussian proxy bounds the variance).
\end{remark}

\begin{theorem}[Error probability of Adjacent-BAI, Theorem 2]
  \label{thm:upper_bound}
  \lean{MaynardZhang2026Complexity.prob_ne_bestArm_le}
  \leanok
  \uses{lem:ls_subgaussian, lem:adjacency_ge, lem:min_gap_best_arm, def:adjacent_bai}
  Let $\norm{x} \le 1$ for $x \in \X$, let the allocation satisfy the rounding guarantee, let
  $\norm{\theta_t} \le 1$ for $t < T$, let $\theta$ have min-gap at least $\Delta > 0$ and best arm
  $x_\star$, and let the noise be centered $1$-sub-Gaussian. In every run of Adjacent-BAI,
  \[
    \Prob(\hat x \ne x_\star) \le \abs{\I^{x_\star}}
      \exp\left(-\frac{T}{36\, \Hadj(\X, \Delta)}\right).
  \]
\end{theorem}

\begin{proof}
  \uses{lem:ls_subgaussian, lem:adjacency_ge, lem:min_gap_best_arm, lem:local_global}
  By Lemma~\ref{lem:min_gap_best_arm}, $x_\star$ is the arm of the min-gap, a vertex, and
  $(x_\star - z)^\top \thetabar_T \ge \Delta$ for every $z \in \I^{x_\star}$. If $\hat x \ne x_\star$,
  then $(\hat x - x_\star)^\top \thetahat_T \ge 0$, and Lemma~\ref{lem:adjacency_ge} gives
  $z \in \I^{x_\star}$ with $(z - x_\star)^\top \thetahat_T \ge 0$, hence
  $(z - x_\star)^\top (\thetahat_T - \thetabar_T) \ge \Delta$. By the union bound and the
  sub-Gaussian tail bound (Lemma~\ref{lem:ls_subgaussian}),
  \[
    \Prob(\hat x \ne x_\star) \le \sum_{z \in \I^{x_\star}}
      \exp\left(-\frac{\Delta^2}{18 \norm{z - x_\star}^2_{A^{-1}}}\right),
  \]
  and $\norm{z - x_\star}^2_{A^{-1}} = \frac1T \norm{z - x_\star}^2_{(A/T)^{-1}} \le \frac2T
  \min_\lambda \max_{\I} \norm{x - x'}^2_{A(\lambda)^{-1}}$ by the rounding guarantee. If the value
  of the adjacent design is $0$, there is no adjacent pair, so $\X = \{x_\star\}$
  (Lemma~\ref{lem:local_global}) and the probability is $0$.
\end{proof}

\begin{remark}
  The paper assumes $T \ge d^2$, which serves only to obtain an allocation with the rounding
  guarantee; with the guarantee as a hypothesis it is not needed. The paper's union bound uses the
  strict inequalities $x^\top \thetahat_T > x_\star^\top \thetahat_T$; with the tie-breaking of the
  argmax, the weak form of the adjacency lemma (Lemma~\ref{lem:adjacency_ge}) is used instead.
\end{remark}

\chapter{The stationary fixed-confidence complexity (Section 6 and Appendix D)}
\label{chap:stationary}

For a stationary parameter $\theta$ with unique best arm $x_\star$, the instance-optimal sample
complexity of fixed-confidence best-arm identification is of order
$\min_\lambda \max_{x \in \X \setminus \{x_\star\}} \norm{x_\star - x}^2_{A(\lambda)^{-1}} /
((x_\star - x)^\top \theta)^2$. The paper shows that the maximum can be restricted to the arms
adjacent to $x_\star$.

\begin{lemma}[Ratio at the adjacent arms, Lemma 10]
  \label{lem:ratio_adjacent}
  \lean{MaynardZhang2026Complexity.ratio_le_iSup_adjacentTo}
  \leanok
  \uses{lem:cone_form}
  Let $M$ be positive definite, $\theta \in \R^d$ and $x_\star \in \X$ with
  $y^\top \theta < x_\star^\top \theta$ for all $y \in \X \setminus \{x_\star\}$. For every
  $x \in \X \setminus \{x_\star\}$,
  \[
    \frac{\norm{x_\star - x}_M^2}{((x_\star - x)^\top \theta)^2}
    \le \max_{x' \in \I^{x_\star}} \frac{\norm{x_\star - x'}_M^2}{((x_\star - x')^\top \theta)^2}.
  \]
\end{lemma}

\begin{proof}
  \leanok
  \uses{lem:cone_form}
  With $u_y = (x_\star - y) / (x_\star - y)^\top \theta$, the ratio of $y$ is $\norm{u_y}_M^2$. By
  Lemma~\ref{lem:cone_form}, $u_x$ lies in the convex hull of the $u_z$, $z \in \I^{x_\star}$. The
  function $u \mapsto \norm{u}_M = \norm{M^{1/2} u}$ is convex, so its value at a point of the
  convex hull of a set is at most its value at some point of the set (the paper's Jensen step);
  the maximum is over a finite set.
\end{proof}

\begin{proposition}[Proposition 1]
  \label{prop:stat_equiv}
  \lean{MaynardZhang2026Complexity.iInf_iSup_ratio_eq}
  \leanok
  \uses{lem:ratio_adjacent, def:design, lem:local_global}
  Let $\theta \in \R^d$ and $x_\star \in \X$ with $y^\top \theta < x_\star^\top \theta$ for all
  $y \in \X \setminus \{x_\star\}$. Then
  \[
    \min_\lambda \max_{x \in \X \setminus \{x_\star\}}
      \frac{\norm{x_\star - x}^2_{A(\lambda)^{-1}}}{((x_\star - x)^\top \theta)^2}
    = \min_\lambda \max_{x \in \I^{x_\star}}
      \frac{\norm{x_\star - x}^2_{A(\lambda)^{-1}}}{((x_\star - x)^\top \theta)^2},
  \]
  the minima being over the designs with positive definite design matrix.
\end{proposition}

\begin{proof}
  \leanok
  \uses{lem:ratio_adjacent, lem:vertices_subset, lem:local_global, lem:strict_expose}
  For each design the two maxima are equal. If $\X = \{x_\star\}$, both are maxima over the empty
  set ($0$ in Lean). Otherwise $x_\star$ is a vertex (Lemma~\ref{lem:strict_expose}) with an adjacent
  vertex (Lemma~\ref{lem:local_global}), both maxima are over nonempty finite sets, $\le$ holds by
  Lemma~\ref{lem:ratio_adjacent} and $\ge$ since $\I^{x_\star} \subseteq \X \setminus \{x_\star\}$
  (Lemma~\ref{lem:vertices_subset}). The paper's assumption that $\X$ spans $\R^d$ is not needed.
\end{proof}


\part{Prerequisites to be added to Mathlib and LML}
\label{part:prereq}

\chapter{Vertices and edges of the convex hull of a finite set}
\label{chap:prereq_polytope}

Throughout, $\X$ is a finite subset of a real inner product space of finite dimension, $\V$ its set
of vertices and $\I$ its set of adjacent pairs (Definition~\ref{def:adjacency}).

\begin{lemma}[Vertices are points of the set]
  \label{lem:vertices_subset}
  \lean{Set.vertices_subset, Set.Finite.finite_adjacentTo}
  \leanok
  \uses{def:adjacency}
  $\V \subseteq \X$; in particular $\V$ and $\I$ are finite.
\end{lemma}

\begin{proof}
  \leanok
  An exposed point is an extreme point, and the extreme points of $\conv \X$ lie in $\X$.
\end{proof}

\begin{lemma}[Maximizers over a convex hull]
  \label{lem:max_face}
  \lean{Set.mem_convexHull_setOf_inner_eq}
  \leanok
  Let $S$ be any subset of a real inner product space, $w$ a vector and $c \in \R$ with
  $y^\top w \le c$ for every $y \in S$. If $z \in \conv S$ and $z^\top w \ge c$, then $z$ is a
  convex combination of the points $y \in S$ with $y^\top w = c$.
\end{lemma}

\begin{proof}
  \leanok
  Split $S = S_= \cup S_<$ according to whether $y^\top w = c$ or $y^\top w < c$. If $S_<$ is
  empty there is nothing to prove; if $S_=$ is empty, $\conv S$ lies in the convex set
  $\{u : u^\top w < c\}$, a contradiction. Otherwise $\conv S$ is the union of the segments
  $[a, b]$ with $a \in \conv S_=$ and $b \in \conv S_<$; writing $z = (1 - t) a + t b$, we get
  $c \le z^\top w = (1 - t) c + t\, b^\top w$ with $b^\top w < c$, so $t = 0$ and
  $z = a \in \conv S_=$.
\end{proof}

\begin{lemma}[Vertices of a finite set]
  \label{lem:strict_expose}
  \lean{Set.mem_vertices_iff_exists_inner_lt, Set.mem_vertices_of_inner_lt,
    Set.exists_inner_lt_of_mem_vertices, Set.Finite.mem_vertices_of_mem_extremePoints,
    Set.Finite.convexHull_vertices}
  \leanok
  \uses{lem:vertices_subset, lem:max_face}
  For $x \in \X$, the following are equivalent: $x$ is a vertex; there is $w$ with
  $y^\top w < x^\top w$ for every $y \in \X \setminus \{x\}$; $x$ is an extreme point of
  $\conv \X$. Moreover $\conv \X = \conv \V$, and every point of $\X$ is a convex combination of
  vertices.
\end{lemma}

\begin{proof}
  \leanok
  \uses{lem:vertices_subset, lem:max_face}
  If $x$ is exposed by $w$, it is the unique maximizer of $w$ over $\conv \X$, hence strictly better
  than the other points of $\X$. If $x$ strictly beats $\X \setminus \{x\}$ for $w$, every
  maximizer of $w$ over $\conv \X$ is a convex combination of the maximizers in $\X$
  (Lemma~\ref{lem:max_face}), that is of $x$ alone: $x$ is the unique maximizer, an exposed point
  (this direction does not use the finiteness of $\X$). If $x$ is an extreme point of $\conv \X$,
  then $x \in \X$ and $x \notin \conv(\X \setminus \{x\})$ (otherwise $x$ would be an extreme point
  of $\conv(\X \setminus \{x\})$, hence a point of $\X \setminus \{x\}$); separate $x$ strictly from
  this compact convex set (Hahn--Banach) to get $w$. Finally, by the Krein--Milman theorem,
  $\conv \X$ (compact and convex) is the closure of the convex hull of its extreme points, which
  are vertices; as $\V$ is finite, $\conv \V$ is closed, so $\conv \X \subseteq \conv \V$.
\end{proof}

\begin{lemma}[Linear forms are maximized at vertices]
  \label{lem:max_at_vertex}
  \lean{Set.Finite.exists_mem_vertices_forall_inner_le}
  \leanok
  \uses{lem:strict_expose}
  If $\X$ is nonempty, for every $\theta$ there is a vertex $v$ with $y^\top \theta \le v^\top \theta$
  for every $y \in \X$.
\end{lemma}

\begin{proof}
  \leanok
  \uses{lem:strict_expose}
  $\V$ is finite and nonempty ($\conv \V = \conv \X \ne \emptyset$); take $v \in \V$ maximizing
  $\theta$ over $\V$. The half-space $\{u : u^\top \theta \le v^\top \theta\}$ is convex and
  contains $\V$, hence $\conv \V = \conv \X \supseteq \X$.
\end{proof}

\begin{lemma}[Lexicographic maximizers]
  \label{lem:lex_expose}
  \lean{Set.Finite.exists_inner_lt_of_lex}
  \leanok
  Let $x$ maximize $\varphi$ over $\X$, and let $y^\top \psi < x^\top \psi$ for every
  $y \in \X \setminus \{x\}$ with $y^\top \varphi = x^\top \varphi$. Then for some $\delta > 0$,
  $y^\top (\varphi + \delta \psi) < x^\top (\varphi + \delta \psi)$ for every $y \in \X \setminus \{x\}$.
\end{lemma}

\begin{proof}
  \leanok
  For each of the finitely many $y \in \X \setminus \{x\}$, the inequality holds for all small
  enough $\delta > 0$: for every $\delta > 0$ if $y^\top \varphi = x^\top \varphi$, and by
  continuity in $\delta$ if $y^\top \varphi < x^\top \varphi$.
\end{proof}

\begin{lemma}[Adjacency through a direction]
  \label{lem:is_adjacent_iff}
  \lean{Set.Finite.isAdjacent_iff_exists_inner}
  \leanok
  \uses{lem:strict_expose, lem:max_face}
  Two distinct vertices $x, x'$ are adjacent if and only if there is $w$ with
  $x^\top w = x'^\top w$ and $y^\top w < x^\top w$ for every vertex $y \notin \{x, x'\}$. The only
  vertices in $[x, x']$ are $x$ and $x'$. In particular, if $x$ and $x'$ are adjacent there are
  $w$ and $\eps > 0$ with $x^\top w = x'^\top w \ge y^\top w + \eps$ for every vertex
  $y \notin \{x, x'\}$ (the paper's Eq.~(adjacentdef)).
\end{lemma}

\begin{proof}
  \leanok
  \uses{lem:strict_expose, lem:max_face}
  If $[x, x']$ is the set of maximizers over $\conv \X$ of the linear form $w$, then $x$ and $x'$ are
  both maximizers, and a vertex $y \notin \{x, x'\}$ that is a maximizer lies in $[x, x']$, hence in
  the open segment: impossible for an extreme point of $\conv \X$. Conversely, all vertices have
  $y^\top w \le x^\top w$, hence so does $\conv \V = \conv \X$ (Lemma~\ref{lem:strict_expose}), and
  $[x, x']$ consists of maximizers; a maximizer over $\conv \X = \conv \V$ is a convex combination of
  the maximizing vertices (Lemma~\ref{lem:max_face}), which are $x$ and $x'$, so it lies in
  $[x, x']$. For the last claim take $\eps$ the smallest gap over the finitely many vertices.
\end{proof}

\begin{lemma}[Local optimality implies global optimality]
  \label{lem:local_global}
  \lean{Set.Finite.exists_mem_adjacentTo_inner_pos, Set.Finite.adjacentTo_nonempty}
  \leanok
  \uses{lem:strict_expose, lem:max_at_vertex, lem:lex_expose, lem:is_adjacent_iff}
  Let $x \in \V$.
  \begin{enumerate}
    \item If $(y - x)^\top \theta > 0$ for some $y \in \X$, then $(z - x)^\top \theta > 0$ for some
      $z \in \I^x$.
    \item If $\X$ has a point other than $x$, then $\I^x \ne \emptyset$.
  \end{enumerate}
\end{lemma}

\begin{proof}
  \leanok
  \uses{lem:strict_expose, lem:max_at_vertex, lem:lex_expose, lem:is_adjacent_iff}
  Claim 2 follows from claim 1 with $\theta = y - x$. For claim 1, let $w$ strictly expose $x$
  (Lemma~\ref{lem:strict_expose}). Every $y \in \X \setminus \{x\}$ is $y = x + d_y a_y$ with
  $d_y = (x - y)^\top w > 0$ and $a_y = (y - x)/d_y$, so that $a_y^\top w = -1$.

  By Lemma~\ref{lem:max_at_vertex} applied to the finite set $A = \{a_y : y \in \X \setminus \{x\}\}$,
  some vertex $b$ of $A$ maximizes $\theta$ over $A$; in particular
  $b^\top \theta \ge a_{y}^\top \theta = (y - x)^\top \theta / d_y > 0$ for the given $y$. Let $\psi$
  strictly expose $b$ in $A$, and among the points $y'$ with $a_{y'} = b$ (the points of $\X$ on the
  ray $x + \R_+ b$) let $z$ maximize $d_{y'}$. Set $\varphi = \psi + (b^\top \psi)\, w$. For
  $y' \in \X \setminus \{x\}$,
  \[
    (y' - x)^\top \varphi = d_{y'} \bigl(a_{y'}^\top \psi - b^\top \psi\bigr) \le 0,
  \]
  with equality if and only if $a_{y'} = b$. So $x$ and $z$ maximize $\varphi$ over $\X$, and every
  other maximizer $y'$ lies on the ray with $d_{y'} < d_z$, that is in the open segment $(x, z)$.

  $z$ is a vertex: it maximizes $\varphi$, and among the maximizers of $\varphi$ it is the unique
  maximizer of $-w$ ($(-w)^\top y' = (-w)^\top x + d_{y'}$), so Lemma~\ref{lem:lex_expose} gives a
  direction strictly exposing it. $x$ and $z$ are adjacent by Lemma~\ref{lem:is_adjacent_iff} with
  the direction $\varphi$: a vertex $v \notin \{x, z\}$ maximizing $\varphi$ would lie in the open
  segment $(x, z)$, which is impossible for an extreme point of $\conv \X$. Finally
  $(z - x)^\top \theta = d_z\, b^\top \theta > 0$.
\end{proof}

\begin{lemma}[Adjacency lemma, weak form]
  \label{lem:adjacency_ge}
  \lean{Set.Finite.exists_mem_adjacentTo_inner_nonneg}
  \leanok
  \uses{lem:local_global}
  Let $x \in \V$. If $(y - x)^\top \theta \ge 0$ for some $y \in \X \setminus \{x\}$, then
  $(z - x)^\top \theta \ge 0$ for some $z \in \I^x$.
\end{lemma}

\begin{proof}
  \leanok
  \uses{lem:local_global}
  Otherwise $(z - x)^\top \theta < 0$ for every $z$ in the finite set $\I^x$, hence
  $(z - x)^\top (\theta + \eps (y - x)) < 0$ for every $z \in \I^x$ and every small enough
  $\eps > 0$. But $(y - x)^\top(\theta + \eps (y - x)) \ge \eps \norm{y - x}^2 > 0$, and
  Lemma~\ref{lem:local_global} gives $z \in \I^x$ with $(z - x)^\top (\theta + \eps(y - x)) > 0$.
\end{proof}

\begin{lemma}[Cone property]
  \label{lem:cone_form}
  \lean{Set.Finite.smul_mem_convexHull_image_adjacentTo}
  \leanok
  \uses{lem:local_global}
  Let $\theta \in \R^d$ and $x_\star \in \X$ with $y^\top \theta < x_\star^\top \theta$ for every
  $y \in \X \setminus \{x_\star\}$. For $y \in \X \setminus \{x_\star\}$ let
  $u_y = (x_\star - y) / (x_\star - y)^\top \theta$. Then for every $x \in \X \setminus \{x_\star\}$,
  $u_x \in \conv\{u_z : z \in \I^{x_\star}\}$.
\end{lemma}

\begin{proof}
  \leanok
  \uses{lem:local_global, lem:strict_expose}
  $x_\star$ is a vertex (Lemma~\ref{lem:strict_expose}). The set $K = \conv\{u_z : z \in
  \I^{x_\star}\}$ is compact and convex, and lies with $u_x$ in the affine hyperplane
  $\{u : u^\top \theta = 1\}$. If $u_x \notin K$, a strict separation gives $\phi$ and $c$ with
  $\phi^\top u_x > c > \phi^\top u$ for all $u \in K$; then $\psi = c\,\theta - \phi$ satisfies
  $\psi^\top u_z > 0$ for $z \in \I^{x_\star}$ and $\psi^\top u_x < 0$, that is
  $(z - x_\star)^\top \psi < 0$ for every $z \in \I^{x_\star}$ and $(x - x_\star)^\top \psi > 0$,
  contradicting Lemma~\ref{lem:local_global}. (If $\I^{x_\star} = \emptyset$ the same argument
  applies with $K = \emptyset$.)
\end{proof}

\chapter{Information-theoretic tools}
\label{chap:prereq_information}

These lemmas are in \texttt{LMLPapers} (\texttt{ForMathlib/InformationTheory/KullbackLeibler/},
\texttt{Identification/ChangeOfMeasure.lean}) or in LML
(\texttt{SequentialLearning/DivergenceDecomposition.lean}), and are to be copied or adapted.

\begin{lemma}[Bretagnolle--Huber inequality]
  \label{lem:bretagnolle_huber}
  \lean{InformationTheory.bretagnolle_huber}
  \leanok
  For probability measures $\mu, \nu$ and a measurable set $E$,
  $\mu(E) + \nu(E^c) \ge \frac12 \exp(-\KL(\mu \,\|\, \nu))$.
\end{lemma}

\begin{proof}
  \leanok
  $\mu(E) + \nu(E^c) \ge \int \min(d\mu, d\nu) \ge \frac12 (\int \sqrt{d\mu\, d\nu})^2 \ge \frac12
  \exp(-\KL(\mu \,\|\, \nu))$, by the Cauchy--Schwarz inequality and Jensen's inequality.
\end{proof}

\begin{lemma}[Divergence between Gaussian laws]
  \label{lem:kl_gaussian}
  \lean{InformationTheory.klDiv_gaussianReal_one, InformationTheory.klDiv_gaussianReal_same_var,
    InformationTheory.klDiv_gaussianReal}
  \leanok
  $\KL(\Gauss(m, 1) \,\|\, \Gauss(m', 1)) = (m - m')^2 / 2$. In particular the reward laws of the
  arm $x$ in the linear bandits of parameters $\theta$ and $\theta'$ with standard Gaussian noise
  have divergence $(x^\top(\theta - \theta'))^2 / 2$.
\end{lemma}

\begin{proof}
  \leanok
  Direct computation of the integral of the log-likelihood ratio, an affine function.
\end{proof}

\begin{lemma}[Data processing for the output]
  \label{lem:data_processing_output}
  \lean{Learning.IdentAlg.IsRun.klDiv_map_out_le,
    Learning.IdentAlg.IsFixedBudget.klDiv_outputMeasure_le}
  \leanok
  For an identification algorithm with budget $T$ and runs against two environments, the divergence
  between the laws of the outputs is at most the divergence between the laws of the histories of
  the first $T$ rounds.
\end{lemma}

\begin{proof}
  \leanok
  The output is drawn from the same kernel applied to the history of the $T$ rounds; data
  processing inequality for kernels.
\end{proof}

\begin{lemma}[Chain rule for non-stationary bandits]
  \label{lem:chain_rule_nonstationary}
  \lean{Learning.IsAlgEnvSeq.klDiv_map_history_banditSeq_eq_sum,
    Learning.IsAlgEnvSeq.klDiv_map_history_banditSeq,
    Learning.IsAlgEnvSeq.klDiv_map_history_compProd_banditSeq, Learning.stepKernel_banditSeq}
  \leanok
  For runs of one algorithm against the non-stationary bandits $\nu$ and $\nu'$, the divergence
  between the laws of the histories of the first $T$ rounds is
  $\sum_{t < T} \sum_a \Prob_\nu(X_t = a)\, \KL(\nu_{t, a} \,\|\, \nu'_{t, a})$ (finitely many arms).
\end{lemma}

\begin{proof}
  \leanok
  LML's chain rule for histories (\texttt{IsAlgEnvSeq.klDiv\_map\_history\_stepKernel}): the
  divergence of the step of round $t$ given the history is that of the reward kernels, the policy
  kernels being equal; the divergence of $\Prob_\nu(X_t \in \cdot) \otimes \nu_{t, \cdot}$ and
  $\Prob_\nu(X_t \in \cdot) \otimes \nu'_{t, \cdot}$ is the average of the divergences.
\end{proof}

\chapter{Concentration}
\label{chap:prereq_concentration}

\begin{lemma}[Hoeffding's lemma]
  \label{lem:hoeffding_lemma}
  \lean{ProbabilityTheory.hasSubgaussianMGF_of_mem_Icc_of_integral_eq_zero}
  \leanok
  A centered random variable with values in $[a, b]$ is sub-Gaussian with variance proxy
  $((b - a)/2)^2$.
\end{lemma}

\begin{proof}
  \leanok
  In Mathlib.
\end{proof}

\begin{lemma}[Sampling without replacement]
  \label{lem:sampling_without_replacement}
  \uses{lem:hoeffding_lemma}
  Let $a_0, \dots, a_{n-1}$ and $g_0, \dots, g_{n-1}$ be vectors with $\norm{g_k} \le M$,
  $\bar g = \frac1n \sum_k g_k$, and $\pi$ a uniform permutation of $\{0, \dots, n - 1\}$. Then
  $\sum_s a_s^\top (g_{\pi(s)} - \bar g)$ is sub-Gaussian with variance proxy
  $4M^2 \sum_{s < n - 1} \norm{d_s}^2$, where $d_s = a_s - \frac{1}{n - 1 - s} \sum_{s' > s} a_{s'}$.
\end{lemma}

\begin{proof}
  \uses{lem:hoeffding_lemma}
  By induction on $n$; for $n \le 1$ the variable is $0$. Let $n \ge 2$. A uniform permutation of
  $\{0, \dots, n - 1\}$ is the pair of a uniform first value $\pi(0)$ and an independent uniform
  bijection from the remaining positions to the remaining values (\texttt{Equiv.Perm.decomposeFin}).
  The average of the $g_k$, $k \ne \pi(0)$, is $\bar g' = (n \bar g - g_{\pi(0)}) / (n - 1)$, whence
  \[
    \sum_s a_s^\top (g_{\pi(s)} - \bar g) = d_0^\top (g_{\pi(0)} - \bar g)
      + \sum_{s \ge 1} a_s^\top (g_{\pi(s)} - \bar g').
  \]
  The first term is centered and lies in $[-2M \norm{d_0}, 2M \norm{d_0}]$ (Hoeffding's lemma,
  proxy $4M^2 \norm{d_0}^2$); conditionally on $\pi(0)$, the second is the statistic of the
  remaining positions and values, sub-Gaussian with proxy $4M^2 \sum_{s \ge 1} \norm{d_s}^2$ by the
  induction hypothesis. Multiply the moment generating functions (tower property).
\end{proof}

\begin{lemma}[Orthogonal columns]
  \label{lem:column_bessel}
  For vectors $a_0, \dots, a_{n-1}$ and $d_s = a_s - \frac{1}{n - 1 - s} \sum_{s' > s} a_{s'}$,
  $\sum_{s < n - 1} \norm{d_s}^2 \le 2 \sum_s \norm{a_s}^2$.
\end{lemma}

\begin{proof}
  $d_s = \sum_i D_{i s} a_i$ for the $n \times (n - 1)$ matrix with $D_{ss} = 1$,
  $D_{is} = -1/(n - 1 - s)$ for $i > s$ and $0$ otherwise. Its columns are orthogonal, of squared
  norms $1 + 1/(n - 1 - s) \le 2$. Coordinatewise, for a vector $r \in \R^n$,
  $\sum_s (D_s^\top r)^2 \le 2 \sum_s (\hat D_s^\top r)^2 \le 2 \norm{r}^2$ by Bessel's inequality for
  the orthonormal family of the normalized columns $\hat D_s$; sum over the coordinates.
\end{proof}

\begin{lemma}[Weighted sums of independent sub-Gaussian noises]
  \label{lem:weighted_noise_subgaussian}
  Let $\eps_0, \dots, \eps_{T-1}$ be independent centered $1$-sub-Gaussian variables, independent of
  a random variable $\pi$, and $c_t(\pi)$ real weights. Conditionally on $\pi$,
  $\sum_t c_t(\pi) \eps_t$ is sub-Gaussian with variance proxy $\sum_t c_t(\pi)^2$; in particular
  $\E[e^{\mu S(\pi)} e^{\mu \sum_t c_t(\pi) \eps_t}] \le \E[e^{\mu S(\pi)}]\, e^{\mu^2 C / 2}$
  for any function $S$ when $\sum_t c_t(\pi)^2 \le C$.
\end{lemma}

\begin{proof}
  By independence, integrate first over the noises (Fubini on the product law): for fixed $\pi$ the
  moment generating function of the weighted sum is the product of those of the $c_t \eps_t$.
\end{proof}

